\clearpage
\section{Semi-Supervised Learning as a Solution}

Semi-Supervised Learning (SSL) lies between supervised and unsupervised learning. Its main goal
is to leverage a small labeled dataset along with a large pool of unlabeled data to improve
model generalization, especially when acquiring labels is expensive or time-consuming. This
makes SSL particularly relevant for materials discovery, where every label costs a DFT
calculation or an experiment (Section~1).

In this project, we explore five SSL algorithms. The wrapper methods are built on
\textbf{tree-ensemble classifiers}: Random Forest as the reference learner (as in classical
QSAR/QSPR studies) and a LightGBM gradient-boosting learner tuned with Optuna:

\begin{itemize}
  \item \textbf{Self-Learning (Pseudo-Labeling):} a tree-ensemble model is trained on labeled
  data, then used to predict labels for unlabeled data. High-confidence predictions are added
  iteratively to the training set to refine the model.
  \item \textbf{Co-Training (Multi-View Learning):} two models are trained on different
  feature subsets (views) of the same data: elemental-property statistics and
  stoichiometry/valence-electron descriptors. Each model labels unlabeled data for the other,
  exploiting complementary information.
  \item \textbf{Label Propagation:} models the dataset as a graph where nodes are materials and
  edges represent compositional similarity. Labels propagate from labeled to unlabeled nodes,
  respecting the data manifold.
  \item \textbf{Label Spreading:} similar to label propagation, but includes a normalization
  step to ensure a smooth label distribution across the graph and to reduce sensitivity to
  noisy edges.
  \item \textbf{Cluster-then-Label:} an unsupervised K-Means clustering of labeled and
  unlabeled materials; unlabeled members of label-pure clusters receive the cluster's label
  (the \emph{cluster assumption} of SSL).
\end{itemize}

\begin{table}[H]
\centering\small
\renewcommand{\arraystretch}{1.25}
\begin{tabularx}{\textwidth}{>{\bfseries}l X X X X X}
\toprule
Feature & Self-Learning & Co-Training & Label Propagation & Label Spreading & Cluster-then-Label\\
\midrule
Approach & Iterative pseudo-labeling & Multi-view learning & Graph-based, hard labels &
Graph-based, soft labels & Clustering + label transfer\\
Strengths & Simple, easy to implement, scalable & Exploits complementary features, reduces bias
& Captures manifold structure, smooth predictions & Smooths labels, robust to noise & Very fast,
uses global structure of the pool\\
Limitations & Sensitive to initial model errors & Requires distinct views, can be complex &
Computationally intensive, sensitive to graph quality & Same as label propagation, can
over-smooth & Fails when clusters mix classes\\
\bottomrule
\end{tabularx}
\end{table}

By comparing these SSL methods on our materials datasets, we aim to identify which technique
best balances prediction accuracy and robustness when labeled data is limited.
